\begin{definition}[Transitive set]
\label{def:transitive-set}
A set \(A\) is \emph{transitive} iff
\[
  \forall x\in A\,\forall y\in x\,(y\in A),
\]
equivalently, iff \(\bigcup A\subseteq A\).
\end{definition}

\begin{definition}[Ordinal and cardinal]
\label{def:ordinal-cardinal}
An \emph{ordinal} is a transitive set \(\alpha\) such that \(\in\) strictly
well-orders \(\alpha\).  For ordinals \(\alpha,\beta\), write
\(\alpha<\beta\) iff \(\alpha\in\beta\).  A \emph{cardinal} is an ordinal
\(\kappa\) for which no smaller ordinal is bijective with \(\kappa\).  Thus a
cardinal is already the least ordinal of its bijection type; the phrase
``initial ordinal'' will not be used as a second term.  A cardinal is
\emph{infinite} iff it is not bijective with any \(n\in\N\).
\end{definition}

For an ordinal \(\alpha\), its \emph{successor ordinal} is
\(\alpha+1\defeq\alpha\cup\{\alpha\}\).  A \emph{limit ordinal} is a nonzero
ordinal that is not a successor ordinal.  These terms classify ordinals by
their immediate-predecessor behavior; they do not introduce a second order or
a second notion of cardinality.  Limit ordinals index the elementary chains
used later, and successor and limit stages distinguish the two clauses of the
Henkin recursion.

\begin{theorem}[Transfinite recursion]
\label{thm:transfinite-recursion}
Let \(\delta\) be an ordinal.  Suppose a functional rule \(G\) assigns a set
\(G(g)\) to every function \(g\) whose domain is an ordinal smaller than
\(\delta\).  Then there is a unique function \(F\) with domain \(\delta\) such
that
\[
 F(\alpha)=G(F\mathbin{\upharpoonright}\alpha)
 \qquad\text{for every }\alpha<\delta.
\]
\textit{Status: cited theorem.}
\end{theorem}

See \cite{Jech2003} for a complete proof.  The theorem is the exact
metalinguistic recursion principle used below; it is not a new inference rule
of any object-language calculus.

